% 柯伊伯带（综述）
% license CCBYSA3
% type Wiki

本文根据 CC-BY-SA 协议转载翻译自维基百科\href{https://en.wikipedia.org/wiki/Kuiper_belt}{相关文章}。

柯伊伯带（/ˈkaɪpər/ ⓘ）$^{[2]}$ 是位于太阳系外部的一种环星盘结构，其范围从海王星轨道处约 30 个天文单位（AU）一直延伸到距离太阳约 50 AU 的区域。$^{[3]}$ 它在结构上类似于小行星带，但规模要大得多——宽度约为其 20 倍，总质量约为其 20–200 倍。$^{[4][5]}$ 与小行星带相似，柯伊伯带主要由太阳系形成早期遗留下来的小天体或残骸组成。不同之处在于，许多小行星主要由岩石和金属构成，而大多数柯伊伯带天体则主要由冻结的挥发性物质（通常称为“冰”）构成，例如甲烷、氨和水冰。柯伊伯带是天文学界普遍认可的多颗矮行星的家园，包括 Orcus、冥王星（Pluto）$^{[6]}$、妊神星（Haumea）$^{[7]}$、夸奥瓦（Quaoar）以及鸟神星（Makemake）$^{[8]}$。太阳系中的一些卫星，例如海王星的海卫一（Triton）和土星的菲比（Phoebe），也可能起源于这一地区。$^{[9][10]}$

柯伊伯带以荷兰天文学家杰拉德·柯伊伯（Gerard Kuiper）的名字命名，他在 1951 年提出了这种带状结构存在的一种设想。$^{[11]}$ 在他之前和之后，也有研究者提出过类似的假说，例如 20 世纪 30 年代的肯尼思·埃奇沃斯（Kenneth Edgeworth）。$^{[12]}$ 对该结构最为直接的预测来自天文学家胡利奥·安赫尔·费尔南德斯（Julio Ángel Fernández），他在 1980 年发表论文，提出在海王星之外存在一条彗星带$^{[13][14]}$，并认为它可能是短周期彗星的来源。$^{[15][16]}$

1992 年，小行星 15760 Albion 被发现，这是继 1930 年发现冥王星、1978 年发现卡戎之后确认的首个柯伊伯带天体（KBO）。$^{[17]}$ 自其发现以来，已知的柯伊伯带天体数量已增加到数千个，并且人们认为直径超过 100 km（62 英里）的 KBO 可能多达 100000 个以上。$^{[18]}$ 起初，柯伊伯带被认为是周期彗星（轨道周期小于 200 年）的主要储存库。然而，自 20 世纪 90 年代中期以来的研究表明，柯伊伯带在动力学上是稳定的，而彗星真正的起源地是散射盘——这是一个由海王星在约 45 亿年前向外迁移所形成的动力学活跃区域。$^{[19]}$ 散射盘天体（例如阋神星 Eris）具有极高的轨道偏心率，其轨道最远可延伸至距离太阳约 100 AU 的位置。$^{[a]}$

柯伊伯带不同于假想中的奥尔特云，后者被认为距离太阳远得多（约为其一千倍），且整体上近似球形。柯伊伯带内的天体，与散射盘成员以及任何潜在的希尔斯云或奥尔特云天体，统称为海王星外天体（TNO）。$^{[22]}$ 冥王星是目前已知柯伊伯带中最大、质量最大的成员，同时也是已知 TNO 中体积最大、质量第二大的天体，仅次于散射盘中的阋神星。$^{[a]}$ 冥王星最初被视为一颗行星，但由于其属于柯伊伯带成员，这一事实促使其在 2006 年被重新归类为矮行星。冥王星在成分上与许多其他柯伊伯带天体相似，其轨道周期也具有一类 KBO 的典型特征，这类天体被称为“冥族天体”（plutinos），它们与海王星处于相同的 2:3 轨道共振状态。

柯伊伯带与海王星有时也被视为太阳系范围的标志之一，其他替代界限还包括日球层顶（heliopause），以及太阳引力影响与其他恒星引力影响相当的距离（估计约在 50000–125000 AU 之间）。$^{[23]}$
\subsection{历史}
\begin{figure}[ht]
\centering
\includegraphics[width=6cm]{./figures/c3c318a9ff09ac52.png}
\caption{冥王星与卡戎} \label{fig_Kuiper_1}
\end{figure}
1930 年冥王星被发现之后，许多人开始推测它可能并非孤立存在。如今被称为柯伊伯带的这一区域，在随后数十年间以多种不同形式被提出为假想结构。直到 1992 年，人们才首次获得其存在的直接证据。由于在此之前关于柯伊伯带性质的设想在数量和形式上都十分多样，至今仍难以确定究竟应将最早提出这一概念的功劳归于谁。$^{[24]:106}$
\subsubsection{假说}
\begin{figure}[ht]
\centering
\includegraphics[width=6cm]{./figures/acd1d13b7d423c69.png}
\caption{以其命名柯伊伯带的天文学家杰拉德·柯伊伯} \label{fig_Kuiper_2}
\end{figure}

\subsection{另见}
\begin{itemize}
\item 小行星带
\item 可能的矮行星列表
\item 海王星外天体列表
\item 第九行星
\item 海王星外天体光谱类型
\end{itemize}
\subsection{注释}
a.关于术语“散射盘”和“柯伊伯带”的使用，文献中存在不一致之处。对一些人而言，这两个概念是截然不同的天体群；而另一些人则认为，散射盘是柯伊伯带的一部分。甚至有作者在同一篇出版物中会在这两种用法之间交替使用。由于国际天文学联合会下属的小行星中心——负责编目太阳系内小行星的机构——做出了这一区分，维基百科关于跨海王星区域的文章在编辑时也采用了同样的区分。在维基百科上，已知质量最大的海王星外天体厄里斯并不属于柯伊伯带，这使得冥王星成为质量最大的柯伊伯带天体。
\subsection{参考文献}
\begin{enumerate}
\item 来源：小行星中心, www.cfeps.net等。
\item “柯伊伯带”.Dictionary.com 精编版（在线）。无日期。；Merriam-Webster.com 词典。梅里亚姆-韦伯斯特。“发音提示：快速参考”Merriam-Webster.com梅里亚姆-韦伯斯特已归档自原始网页于2007年6月1日存档。 检索于2025年4月3日
\item 斯特恩，艾伦；科尔韦尔，乔舒亚·E.（1997）。“原始埃奇沃斯-柯伊伯带中的碰撞侵蚀与30至50天文单位柯伊伯空隙的形成”. 天体物理学报. 490 (2): 879–882. 文献标识码:1997ApJ...490..879S. doi:10.1086/304912.
\item 德尔桑蒂，奥黛丽 & 朱维特，大卫（2006）。行星之外的太阳系（PDF）。夏威夷大学天文学研究所。文献标识码:2006ssu..book..267D。于2007年9月25日从原文存档。 检索于2007年3月9日。
\item 克拉斯金斯基，G. A.; 彼捷娃，E. V.；瓦西里耶夫，M. V.；亚古季娜，E. I. （2002年7月）。“小行星带中的暗物质”。国际行星科学杂志. 158 (1): 98–105. 天文学数据库标识符:2002Icar..158...98K. DOI:10.1006/icar.2002.6837.
\item 克里斯滕森，拉斯·林德伯格。“国际天文学联合会2006年大会：国际天文学联合会决议投票结果”。国际天文学联合会。于2014年4月29日存档自原始网页。检索日期：2021年5月25日。
\item 克里斯滕森，拉斯·林德伯格（2008年9月17日）。“国际天文学联合会命名第五颗矮行星为妊神星”。国际天文学联合会。于2014年4月25日存档自原文。检索日期：2021年5月25日。
\item 克里斯滕森，拉斯·林德伯格（2008年7月19日）。“第四颗矮行星被命名为马凯马克”。国际天文学联合会。于2019年6月16日存档自原文。检索日期：2021年5月25日。
\item 约翰逊，托伦斯·V.；以及卢宁，乔纳森·I.土星的卫星菲比作为来自太阳系外层的捕获天体, 自然，第435卷，第69–71页。
\item 克雷格·B·阿格诺尔与道格拉斯·P·汉密尔顿（2006）。“海王星在双行星引力遭遇中捕获其卫星海卫一”（PDF）。自然441（7090）：192–194。天文学编码：2006Natur.441..192ADOI10.1038/nature04792PMID16688170S2CID4420518。于2007年6月21日从原始网页存档。检索日期：2006年6月20日
\item 开普勒，G. P.（1951）。“关于太阳系的起源”。载于海内克，J. A.（编）。天体物理学：专题研讨会。美国纽约市，纽约州：麦格劳-希尔出版社。第357–424.
\item “柯伊伯带：事实——美国国家航空航天局科学”。2017年11月14日。于2024年3月12日存档自原始网页。检索日期：2024年2月18日。
\item J. A. 费尔南德斯（1980）。“关于海王星外彗星带的存在”. 马德里国家天文台. 192 (3): 481–491. Bibcode:1980MNRAS.192..481F. doi:10.1093/mnras/192.3.481.
\item 莫尔比戴利，A.；托马斯，F.；穆恩斯，M.（1995年12月1日）。“柯伊伯带的共振结构与前五个海王星外天体的动力学”. 《伊卡洛斯》. 118 (2): 322–340. 文献标识码:1995Icar..118..322M. DOI:10.1006/icar.1995.1194. ISSN 0019-1035.
 gunjan.sogani（2022年9月10日）。“柯伊伯带及其成员的发现”。Wondrium每日。于2023年8月1日从原文存档。检索日期：2023年8月1日
 “胡里奥·A·费尔南德斯”.美国国家科学院官网于2023年8月1日存档自原始网页。检索日期：2023年8月1日
 朱维特，大卫；刘佳（1993）。“柯伊伯带候选天体1992 QB1的发现”。自然.362(6422)：730–732。Bibcode：1993Natur.362..730Jdoi10.1038/362730a0S2CID4359389
 “首席研究员的观点”.新视野。2012年8月24日。于2014年11月13日从原始网页存档.
 莱维森，哈罗德·F.；多尼斯，卢克（2007）。“彗星种群与彗星动力学”。载于露西·安·亚当斯·麦克法登；保罗·罗伯特·魏斯曼；托伦斯·V·约翰逊（编）。《太阳系百科全书》（第2版）。阿姆斯特丹；波士顿：学术出版社。第575–588页.ISBN978-0-12-088589-3
 魏斯曼与约翰逊，2007年，《太阳系百科全书》，脚注第584页。
 “半人马星与散盘天体列表”.国际天文学联合会：小行星中心. 哈佛-史密森天体物理中心中央天文电报局。2011年1月3日。已归档于2017年6月29日自原始网页存档。检索于2011年1月3日
 热拉尔·福尔（2004）。“截至2004年5月20日的小行星系统描述”。于2007年5月29日从原文存档。检索日期：2007年6月1日。
 “太阳系的边缘在哪里？”.戈达德媒体工作室. 美国国家航空航天局戈达德太空飞行中心。2017年9月5日。已存档自原始网页于2021年12月16日存档于2019年9月22日
 兰德尔，丽莎（2015）。暗物质与恐龙。纽约：埃科/哈珀柯林斯出版社。ISBN 978-0-06-232847-2.
 “‘柯伊伯带’这一术语有何不妥之处？（或者，为何要以一位不相信其存在的人的名字来命名这一天体？）”。国际彗星季刊于2019年10月8日存档自原始网页。检索日期：2010年10月24日
 戴维斯，约翰·K.；麦克法兰，J.；贝利，马克·E.；马斯登，布赖恩·G.；叶伟文（2008）。“关于跨海王星区域的早期思想发展”（PDF）。载于M. 安东内塔·巴拉奇；赫尔曼·博恩哈特；戴尔·克鲁伊昌克；亚历山德罗·莫尔比德利（编）。
 戴维斯，约翰·K.（2001）。超越冥王星：探索太阳系的外缘。剑桥大学出版社。
 开普勒，杰拉德（1951）。“论太阳系的起源”. 美国国家科学院院刊. 37 (1): 233. doi:10.1073/pnas.37.4.233. PMC 1063291. PMID 16588984.
 大卫·朱维特。“为什么叫‘柯伊伯’带？”。夏威夷大学于2019年2月12日从原始网页存档。检索日期：2007年6月14日
 饶，M. M.（1964）。“向量测度的分解”（PDF）。美国国家科学院院刊51（5）：771–774。文献标识码：1964PNAS...51..771RDOI10.1073/pnas.51.5.771PMC300359PMID16591174已归档自原始网页于2016年6月3日。检索于2007年6月20日
 C. T. 科瓦尔；W. 利勒；B. G. 马尔斯登（1977）。“/2060/ 凯龙的发现与轨道”。载于：太阳系动力学；研讨会论文集. 81。哈雷天文台，哈佛—史密森天体物理中心：245页。文献编码:1979IAUS...81..245K.
 J. V. 斯科蒂；D. L. 拉比诺维茨；C. S. 肖梅克；E. M. 肖梅克；D. H. 利维；T. M. 金；E. F. 海林；J. 阿卢；K. 劳伦斯；R. H. 麦克诺特；L. 弗雷德里克；D. 托伦；B. E. A. 穆勒（1992）。“1992 AD”。国际天文学联合会通告 (5434): 1. 文献标识码:1992IAUC.5434....1S.
 霍纳，J.；埃文斯，N. W.；贝利，M. E.（2004）。“半人马星群的模拟Ⅰ：总体统计”. 皇家天文学会月报. 354 (3): 798–810. arXiv:astro-ph/0407400. 文献编码:2004MNRAS.354..798H. DOI:10.1111/j.1365-2966.2004.08240.x. S2CID 16002759.
 大卫·朱维特（2002）。“从柯伊伯带天体到彗星核：缺失的远红外观测物质”. 《天文期刊》. 123 (2): 1039–1049. 文献标识码:2002AJ....123.1039J. DOI:10.1086/338692. 跨学科索引 122240711.
 奥尔特，J. H.（1950）。“环绕太阳系的彗星云的结构及其起源假说”。荷兰天文研究所公报. 11: 91. 文献标识码:1950BAN....11...91O.
 “柯伊伯带”。九大行星。2019年10月8日。于2021年5月16日存档自原始网页。检索日期：2020年8月16日
 J. A. 费尔南德斯（1980）。“关于海王星外侧存在彗星带的探讨”. 皇家天文学会月报. 192 (3): 481–491. 文献标识码:1980MNRAS.192..481F. DOI:10.1093/mnras/192.3.481.
 邓肯；奎因；特雷梅恩（1988）。“短周期彗星的起源”. 天体物理学报. 328：L69。文献标识码:1988ApJ...328L..69D. doi:10.1086/185162.
 马斯登，B.S.；朱维特，D.；马斯登，B.G.（1993）。“1993 FW”。国际天文学联合会通告（5730）。小行星中心：1。文献标识码:1993IAUC.5730....1L.
 戴奇斯，普雷斯顿（2018年12月14日）。“关于柯伊伯带的10个须知”。美国国家航空航天局太阳系探索于2019年1月10日存档自原始网页。检索日期：2019年12月1日
 “柯伊伯带20岁”。天体生物学杂志。2012年9月1日。于2020年10月30日从原文存档。检索日期：2019年12月1日
 沃森，保罗（2019年1月1日）。“成功穿越冥王星之外的区域，NASA探测器开始传回柯伊伯带天体的图像”。科学。美国科学促进会。于2022年10月8日存档自原始网页。检索日期：2019年12月1日
 克莱德·汤博，《最后的发言》，致编辑的信件，《天空与望远镜》, 1994年12月，第8页
 “‘柯伊伯带’这一术语有何不妥之处？”.国际彗星季刊于2019年10月8日存档自原始网页。检索日期：2021年12月19日
 M. C. 德·桑克蒂斯；M. T. 卡普里亚与A. 科拉迪尼（2001）。“柯伊伯带天体的热演化与分异”. 《天文学杂志》. 121 (5): 2792–2799. 文献标识码:2001AJ....121.2792D. DOI:10.1086/320385.
 “探索太阳系的边缘”。美国科学家网。2003年。于2009年3月15日从原文存档。检索日期：2007年6月23日
 迈克尔·E·布朗；玛格丽特·潘（2004）。“柯伊伯带的平面” （PDF）. 《天文学杂志》. 127 (4): 2418–2423. 文献标识码:2004AJ....127.2418B. DOI:10.1086/382515. S2CID 10263724。于2020年4月12日从原始内容存档 （PDF）。
 佩蒂，让-马克；莫尔比代利，亚历山德罗；瓦尔塞基，乔瓦尼·B.（1998）。“大质量散射行星体与小天体带的激发”（PDF）.《伊卡洛斯》141（2）：367。文献标识码:1999Icar..141..367PDOI10.1006/icar.1999.6166. 于2007年8月9日从原始网页存档。检索日期：2007年6月23日
 卢尼恩，乔纳森·I.（2003）。“柯伊伯带”（PDF）。于2007年8月9日从原文存档。检索日期：2007年6月23日。
 朱维特，D.（2000年2月）。“经典柯伊伯带天体（CKBOs）”。于2007年6月9日从原始网页存档。检索日期：2007年6月23日。”
 默尔丁，P.（2000）。“柯伊伯带天体”。天文学与天体物理学百科全书. Bibcode:2000eaa..bookE5403.. doi:10.1888/0333750888/5403. ISBN 978-0-333-75088-9.
 埃利奥特，J. L.；等（2005）。“深黄道巡天：柯伊伯带天体与半人马星的搜寻。II. 动力学分类、柯伊伯带平面及核心群体”（PDF）。《天文学杂志》129（2）：1117–1162。文献编码:2005AJ....129.1117EDOI10.1086/427395存档于原始网页2013年7月21日检索日期2012年8月18日
 “天体命名：小行星”。国际天文学联合会于2008年12月16日存档自原始网页。检索日期：2008年11月17日
 佩蒂，J.-M.；格拉德曼，B.；卡韦拉尔斯，J.J.；琼斯，R.L.；帕克，J.（2011）。“经典柯伊伯带内核的现实与起源”（PDF）。欧洲行星科学大会—美国天文学会行星科学分会联合会议（2011年10月2日至7日）。文献标识码:2011epsc.conf..722P已于2016年3月4日存档自原始网页。检索于2016年5月4日
 莱维森，哈罗德·F.；莫尔比戴利，亚历山德罗（2003）。“海王星迁移期间天体向外迁移导致柯伊伯带的形成”。自然. 426 (6965): 419–421. 文献标识码:2003Natur.426..419L. DOI:10.1038/nature02120. PMID 14647375. S2CID 4395099.
 斯蒂芬斯，丹尼斯·C.；诺尔，基思·S.（2006）。“利用NICMOS探测六对海王星外双星：冷经典盘中双星比例极高”。天文学杂志. 130 (2): 1142–1148. arXiv:astro-ph/0510130. 文献编码:2006AJ....131.1142S. DOI:10.1086/498715. S2CID 204935715.
 弗雷泽，韦斯利·C.；布朗，迈克尔·E.；莫尔比迪利，亚历山德罗；帕克，亚历克斯；巴蒂金，康斯坦丁（2014）。“柯伊伯带天体的绝对星等分布”。天体物理学报. 782 (2): 100. arXiv:1401.2157. 文献编码:2014ApJ...782..100F. doi:10.1088/0004-637X/782/2/100. S2CID 2410254.
 莱维森，哈罗德·F.；斯特恩，S·艾伦（2001）。“主柯伊伯带倾角分布的尺寸依赖性研究”。天文学杂志. 121 (3): 1730–1735. arXiv:astro-ph/0011325. Bibcode:2001AJ....121.1730L. doi:10.1086/319420. S2CID 14671420.
\end{enumerate}